\documentclass[11pt,a4paper]{article}
% preamble.tex — estilo común de todos los apuntes Froth.
% Cada apunte hace % preamble.tex — estilo común de todos los apuntes Froth.
% Cada apunte hace % preamble.tex — estilo común de todos los apuntes Froth.
% Cada apunte hace \input{preamble} justo después de \documentclass.
\usepackage[margin=2.4cm]{geometry}
\usepackage{xcolor}
\definecolor{litblue}{RGB}{31,79,121}
\definecolor{litgreen}{RGB}{27,94,32}
\definecolor{litorange}{RGB}{230,81,0}
\usepackage{parskip}
\usepackage{titlesec}
\titleformat{\section}{\Large\bfseries\color{litblue}}{\thesection}{0.6em}{}
\titleformat{\subsection}{\large\bfseries\color{litblue}}{\thesubsection}{0.6em}{}
\usepackage{tcolorbox}
\usepackage{fancyhdr}
\pagestyle{fancy}
\fancyhf{}
\fancyhead[L]{\small\color{litblue}Froth — Apuntes de ML}
\fancyhead[R]{\small\thepage}
\renewcommand{\headrulewidth}{0.4pt}
\usepackage[colorlinks=true,linkcolor=litblue,urlcolor=litblue]{hyperref}

% Cabecera de cada apunte: \cabecera{numero}{titulo}
\newcommand{\cabecera}[2]{%
  \begin{tcolorbox}[colback=litblue,coltext=white,colframe=litblue,arc=2mm]
    {\Large\bfseries Apunte #1 — #2}\\[3pt]
    {\small Proyecto Froth · aprender Machine Learning construyendo}
  \end{tcolorbox}\vspace{0.8em}}

% Cajas temáticas
\newtcolorbox{concepto}[1]{colback=blue!4,colframe=litblue,title={Concepto: #1},arc=1.5mm}
\newtcolorbox{practica}{colback=green!5,colframe=litgreen,title={Pruébalo tú (teclado en mano)},arc=1.5mm}
\newtcolorbox{ojo}{colback=orange!8,colframe=litorange,title={Ojo / error típico},arc=1.5mm}
\newtcolorbox{resumen}{colback=gray!8,colframe=gray!60,title={En una frase},arc=1.5mm}

\newcommand{\codigo}[1]{\texttt{#1}}
 justo después de \documentclass.
\usepackage[margin=2.4cm]{geometry}
\usepackage{xcolor}
\definecolor{litblue}{RGB}{31,79,121}
\definecolor{litgreen}{RGB}{27,94,32}
\definecolor{litorange}{RGB}{230,81,0}
\usepackage{parskip}
\usepackage{titlesec}
\titleformat{\section}{\Large\bfseries\color{litblue}}{\thesection}{0.6em}{}
\titleformat{\subsection}{\large\bfseries\color{litblue}}{\thesubsection}{0.6em}{}
\usepackage{tcolorbox}
\usepackage{fancyhdr}
\pagestyle{fancy}
\fancyhf{}
\fancyhead[L]{\small\color{litblue}Froth — Apuntes de ML}
\fancyhead[R]{\small\thepage}
\renewcommand{\headrulewidth}{0.4pt}
\usepackage[colorlinks=true,linkcolor=litblue,urlcolor=litblue]{hyperref}

% Cabecera de cada apunte: \cabecera{numero}{titulo}
\newcommand{\cabecera}[2]{%
  \begin{tcolorbox}[colback=litblue,coltext=white,colframe=litblue,arc=2mm]
    {\Large\bfseries Apunte #1 — #2}\\[3pt]
    {\small Proyecto Froth · aprender Machine Learning construyendo}
  \end{tcolorbox}\vspace{0.8em}}

% Cajas temáticas
\newtcolorbox{concepto}[1]{colback=blue!4,colframe=litblue,title={Concepto: #1},arc=1.5mm}
\newtcolorbox{practica}{colback=green!5,colframe=litgreen,title={Pruébalo tú (teclado en mano)},arc=1.5mm}
\newtcolorbox{ojo}{colback=orange!8,colframe=litorange,title={Ojo / error típico},arc=1.5mm}
\newtcolorbox{resumen}{colback=gray!8,colframe=gray!60,title={En una frase},arc=1.5mm}

\newcommand{\codigo}[1]{\texttt{#1}}
 justo después de \documentclass.
\usepackage[margin=2.4cm]{geometry}
\usepackage{xcolor}
\definecolor{litblue}{RGB}{31,79,121}
\definecolor{litgreen}{RGB}{27,94,32}
\definecolor{litorange}{RGB}{230,81,0}
\usepackage{parskip}
\usepackage{titlesec}
\titleformat{\section}{\Large\bfseries\color{litblue}}{\thesection}{0.6em}{}
\titleformat{\subsection}{\large\bfseries\color{litblue}}{\thesubsection}{0.6em}{}
\usepackage{tcolorbox}
\usepackage{fancyhdr}
\pagestyle{fancy}
\fancyhf{}
\fancyhead[L]{\small\color{litblue}Froth — Apuntes de ML}
\fancyhead[R]{\small\thepage}
\renewcommand{\headrulewidth}{0.4pt}
\usepackage[colorlinks=true,linkcolor=litblue,urlcolor=litblue]{hyperref}

% Cabecera de cada apunte: \cabecera{numero}{titulo}
\newcommand{\cabecera}[2]{%
  \begin{tcolorbox}[colback=litblue,coltext=white,colframe=litblue,arc=2mm]
    {\Large\bfseries Apunte #1 — #2}\\[3pt]
    {\small Proyecto Froth · aprender Machine Learning construyendo}
  \end{tcolorbox}\vspace{0.8em}}

% Cajas temáticas
\newtcolorbox{concepto}[1]{colback=blue!4,colframe=litblue,title={Concepto: #1},arc=1.5mm}
\newtcolorbox{practica}{colback=green!5,colframe=litgreen,title={Pruébalo tú (teclado en mano)},arc=1.5mm}
\newtcolorbox{ojo}{colback=orange!8,colframe=litorange,title={Ojo / error típico},arc=1.5mm}
\newtcolorbox{resumen}{colback=gray!8,colframe=gray!60,title={En una frase},arc=1.5mm}

\newcommand{\codigo}[1]{\texttt{#1}}

\begin{document}
\cabecera{23}{DOI, open access y la lista de lectura}

\section{La feature nueva: del ranking a la lectura}

El buscador ya rankeaba papers; ahora además te lleva a \textbf{leerlos}: cada resultado
trae dos enlaces, \textbf{DOI} (la página oficial del paper) y \textbf{open PDF} (el PDF
gratuito, cuando existe). En tu corpus: \textbf{85 de 126 papers (67\%) tienen PDF de acceso
abierto} y 116/126 tienen DOI.

\begin{concepto}{DOI (Digital Object Identifier)}
El \textbf{DOI} es el ``número de cédula'' permanente de un paper: \codigo{10.1016/j...}.
La URL \codigo{https://doi.org/[DOI]} siempre lleva a la página oficial del editor, aunque
la revista cambie de sitio web. Es la forma correcta de citar y de llegar al texto; si el
paper es de pago, ahí usas tus accesos de universidad.
\end{concepto}

\begin{concepto}{open access (OA)}
Un paper es \textbf{open access} cuando existe una copia legalmente gratuita: en la propia
revista, en repositorios (arXiv, HAL...) o en la página de la universidad del autor. OpenAlex
rastrea dónde está esa copia y nos da la URL (\codigo{oa\_url}). Nosotros solo pedimos un
campo más en el \codigo{select} del pull (\codigo{open\_access}) — una línea de código,
y el pipeline regeneró todo lo demás solo.
\end{concepto}

\begin{ojo}
\textbf{Enlazar sí, descargar masivamente no.} Froth te da el ENLACE y tú haces clic: la
lectura ocurre de tu lado, con tus permisos. Descargar PDFs en masa desde editoriales viola
sus términos de uso (riesgo previsto en el plan desde el día 1). Las herramientas serias
(Connected Papers, etc.) hacen exactamente esto mismo: rankear + enlazar.
\end{ojo}

\section{Bonus del día: el vigilante lento de Streamlit}

La app tardaba mucho en arrancar y el log se llenaba de avisos de \codigo{torchvision}.
Diagnóstico: el \emph{file watcher} de Streamlit (auto-recarga al editar código) intenta
examinar los cientos de módulos de \codigo{transformers}, uno por uno. Solución: apagarlo en
\codigo{.streamlit/config.toml} (\codigo{fileWatcherType = "none"}). Lección doble:
(1) un log ruidoso no siempre es un error fatal — hay que leer QUÉ dice; (2) las herramientas
de desarrollo cómodas a veces cuestan rendimiento.

\begin{resumen}
Cada resultado del buscador ahora enlaza al DOI (página oficial) y al PDF open access cuando
existe (67\% de tu corpus). Enlazamos, no descargamos en masa. Y la app arranca más rápido
tras apagar el file watcher de Streamlit.
\end{resumen}

\end{document}
